\documentclass[10pt,a4paper]{article}
\usepackage{fontspec,amsmath,amssymb,graphicx}
\usepackage[margin=18mm]{geometry}
\setmainfont{Times New Roman}
\setlength{\parindent}{0pt}
\setlength{\parskip}{5pt}
\allowdisplaybreaks
\sloppy
\emergencystretch=2em
\begin{document}
\section*{BT\_2\_3\_07 — Lời giải bổ sung}
\subsection*{Kiến thức và trình tự giải}
\textbf{Mục tiêu.} Tận dụng phần cấp hai có dạng bình phương một đạo hàm theo hướng để giảm PDE thành ODE có tham số.

\textbf{Bước 1 — Nhận ra hướng đạo hàm.} Nếu $D=a\partial_x+b\partial_y$ có hệ số hằng, $D^2u=a^2u_{xx}+2ab u_{xy}+b^2u_{yy}$. Với hệ số biến thiên phải dùng


\[
D^2u=a^2u_{xx}+2ab u_{xy}+b^2u_{yy}+(Da)u_x+(Db)u_y.
\]


So sánh cả hạng cấp một trước khi kết luận PDE bằng $D^2u+\lambda Du+\mu u=f$.

\textbf{Bước 2 — Tìm bất biến $s$ và biến tiến $z$.} Giải $Ds=0$, chọn $Dz=k\ne0$. Kiểm tra $J=\partial(s,z)/\partial(x,y)\ne0$. Khi $k$ hằng, $D=k\partial_z$; PDE trở thành ODE theo $z$, trong đó $s$ là tham số. Nếu còn $v_s$, phải giữ lại, chỉ các trường hợp nó triệt tiêu mới giải trực tiếp như ODE.

\textbf{Bước 3 — Giải ODE.} Với $v_{zz}=g(s,z)$, tích phân hai lần cho $v=zF(s)+G(s)+\int_{z_0}^z(z-t)g(s,t)dt$. Với $a v_{zz}+b v_z+c v=g$, giải đa thức $a\lambda^2+b\lambda+c=0$, giữ hai hàm tùy ý theo $s$. Nghiệm kép có thêm hệ số $z$; có thể tìm nghiệm riêng dạng mũ hoặc đa thức.

\textbf{Bước 4 — Dùng dữ liệu đúng biến.} Tính $u_x=s_xv_s+z_xv_z$ và $u_y=s_yv_s+z_yv_z$. Nếu dữ liệu ở $s=0$, chỉ biết $F(0),G(0)$ hoặc đạo hàm của chúng; không suy ra toàn bộ các hàm. So sánh hệ số của $e^z$, $z$, $z^2$ để tìm các hằng số tương thích.

\textbf{Bước 5 — Kiểm tra tính duy nhất và nghiệm cụ thể.} Khi đường dữ liệu gặp tất cả đặc trưng một cách độc lập, giải hệ cho $F,G$. Khi đường dữ liệu là một đặc trưng, mô tả các hàm còn tùy ý và đưa hai lựa chọn rõ ràng. Cuối cùng thế lại phương trình gốc, không chỉ ODE mới.
\subsection*{Lời giải áp dụng cho bài}
\textbf{(a) Phân loại và tọa độ.} $A=1$, $B=-3$, $C=9$, $\delta=0$. Đặt $D=\partial_x-3\partial_y$, $s=3x+y$, $z=x$, $J=-1$. Vì $Ds=0$, $Dz=1$, PDE parabolic trở thành $v_{zz}=e^z$. Đặc trưng $y=-3x+s$.

\textbf{(b) Nghiệm tổng quát.} Tích phân theo $z$ hai lần:
\[u=e^x+xF(3x+y)+G(3x+y).\]

\textbf{(c) Dùng lần lượt hai đường dữ liệu.} Trên $x=0$, $s=y$, nên $1+G(y)=1+\sin y$: $G(s)=\sin s$. Trên $y=0$,
\[e^x+xF(3x)+\sin3x=e^x+\sin3x+x^3\]
cho $F(3x)=x^2$ khi $x\ne0$; tính liên tục xác định cả $F(0)=0$. Vì vậy $F(s)=s^2/9$ và
\[\boxed{u=e^x+\frac{x(3x+y)^2}{9}+\sin(3x+y).}\]
Giữ $s$ cố định thì $v_{zz}=e^z$, nên công thức thỏa PDE. Thế $y=0$ và $x=0$ kiểm tra đúng hai dữ liệu; chúng cùng bằng $1$ tại gốc. Hai đường dữ liệu đã xác định toàn bộ $F,G$, nên nghiệm là duy nhất trong lớp $C^2$ trên $\mathbb R^2$.

\textbf{Chi tiết tích phân theo biến còn lại.} Với $s=3x+y,z=x$, $u_x=3v_s+v_z,u_y=v_s$,
\[u_{xx}=9v_{ss}+6v_{sz}+v_{zz},\quad u_{xy}=3v_{ss}+v_{sz},\quad u_{yy}=v_{ss}.\]
Phần chính $u_{xx}-6u_{xy}+9u_{yy}$ có các hệ số $9-18+9=0$, $6-6=0$, còn $v_{zz}$. Tích phân $v_{zz}=e^z$ lần đầu được $v_z=e^z+F(s)$, lần sau $v=e^z+zF(s)+G(s)$.
\begin{center}\includegraphics[width=.97\linewidth,height=11cm,keepaspectratio]{Do_thi/BT_2_3_07.png}\end{center}
\textit{Đồ thị minh họa; với nghiệm số, xem phương pháp và tham số trong mã Python kèm theo.}\end{document}